\documentclass{article}

\usepackage{neurips_2026}

\usepackage[utf8]{inputenc} % allow utf-8 input
\usepackage[T1]{fontenc}    % use 8-bit T1 fonts
\usepackage{hyperref}       % hyperlinks
\usepackage{url}            % simple URL typesetting
\usepackage{booktabs}       % professional-quality tables
\usepackage{amsfonts}       % blackboard math symbols
\usepackage{nicefrac}       % compact symbols for 1/2, etc.
\usepackage{microtype}      % microtypography
\usepackage{xcolor}         % colors
\usepackage{graphicx}       % images

\title{Model Incrimination Paper}

\author{%
  % Add authors here before camera-ready
}

\begin{document}

\maketitle

\begin{abstract}
% Write your abstract here

\end{abstract}

\section{Introduction}

\section{Approach}

\section{Case Studies}

\subsection{How to Conduct Investigations}

\subsection{Case Study 1}

\subsubsection{Setting}

\subsubsection{Claims About Intent}

\subsubsection{Evidence}

% Add more case studies as needed:
% \subsection{Case Study N}
% \subsubsection{Setting}
% \subsubsection{Claims About Intent}
% \subsubsection{Evidence}

\section{Methodology Takeaways}
% e.g., prompt counterfactuals are confusing

\section{Conclusion}

\begin{ack}
% Add acknowledgments here (hidden during anonymous review)
\end{ack}

\bibliographystyle{plainnat}
\bibliography{references}

\appendix

\newpage
\input{checklist.tex}

\end{document}
